\documentclass[9pt]{extarticle}

\usepackage[utf8]{inputenc}
\usepackage[a4paper, margin=1.2cm]{geometry}
\usepackage[colorlinks=true, urlcolor=blue]{hyperref}
\usepackage{enumitem}
\usepackage{changepage}

\pagestyle{empty}

\setlist[itemize]{label=\textbullet, noitemsep, topsep=0pt}
\setlength{\parindent}{0pt}

\begin{document}

\begin{center}
	{\huge\bfseries Paul Parisot}
	\\[4pt]
	{\Large Ingénieur Logiciel - Blockchain, Go \& Automatisation}
	\\
\end{center}

\par\vspace{-5pt}
\rule{\linewidth}{0.4pt}
\par\vspace{-5pt}

\begin{center}
	Github: \href{https://github.com/paulogarithm}{paulogarithm}
	\\
	\href{mailto:paul.parisot@epitech.eu}{paul.parisot@epitech.eu}
	\\
\end{center}

\par\vspace{-10pt}
\rule{\linewidth}{0.4pt}
\par\vspace{7pt}

{\large\bfseries À propos}
\par\vspace{-8pt}
\rule{\linewidth}{0.4pt}
\par\vspace{5pt}

Ingénieur logiciel avec de solides bases en {\bfseries Go, blockchain et C/C++}, orienté automatisation, systèmes distribués et intégrité des données.
\\
Expérience concrète en {\bfseries smart contracts et blockchain} : contribution à Gno.land (Go, EVM, +10k lignes de code sur des packages critiques) et implémentation from scratch d'algorithmes cryptographiques (AES, RSA).
\\
Un an d'études à Montréal (McGill University, management) et pratique professionnelle de l'anglais (C1). En préparation de l'agrégation d'informatique.
\par\vspace{5pt}

{\large\bfseries Compétences}
\par\vspace{-8pt}
\rule{\linewidth}{0.4pt}
\par\vspace{5pt}

\hspace*{5pt}{\bfseries Maîtrise}: Golang, C/C++, Python, Lua, Assembleur, Shell, SQL
\\
\hspace*{5pt}{\bfseries Blockchain \& Crypto}: smart contracts, EVM (traduction d'opcodes), AES, RSA, intégrité des données
\\
\hspace*{5pt}{\bfseries Outils}: gdb, Valgrind, Git, Docker, Kubernetes, CI/CD, Environnement Linux
\\
\hspace*{5pt}{\bfseries Langues}: Français, Anglais (C1, compétences professionnelles)

\par\vspace{5pt}

{\large\bfseries Expérience Professionnelle}
\par\vspace{-8pt}
\rule{\linewidth}{0.4pt}
\par\vspace{5pt}

{\bfseries Gno.land - Ingénieur Logiciel Junior (Open-Source)}
\begin{itemize}[noitemsep, topsep=0pt]
	\item Développement de 6+ packages core en Gnolang (variante de Go pour smart contracts), incluant {\bfseries printf et une compatibilité avec l'EVM} (traduction d'opcodes en instructions Go).
	\item Détection et correction de {\bfseries 6 bugs critiques} causant des crashs runtime ; {\bfseries ~10k lignes de code} contribuées.
	\item Rédaction de tests exhaustifs sous des politiques strictes de CI/CD et de revue de code.
\end{itemize}

\par\vspace{5pt}

{\bfseries Faurecia / Forvia - Ingénieur Démo Logicielle}
\begin{itemize}[noitemsep, topsep=0pt]
	\item Conception d'un {\bfseries moteur de calcul d'itinéraire GPS} et d'une interface de recherche de POI, intégrés à une application démo présentée au {\bfseries CES 2024}.
	\item Correction d'un ancien code C++ pour le faire communiquer avec les bons drivers matériels (LED) ; scripts Python de pilotage de hardware.
	\item Livraison de fonctionnalités prêtes pour la production {\bfseries en avance sur le planning}.
\end{itemize}

\par\vspace{5pt}

{\bfseries Epitech / IONIS - Assistant Technique \& Formateur}
\begin{itemize}[noitemsep, topsep=0pt]
	\item Enseignement à {\bfseries plus de 400 étudiants} en C, C++, asm et Haskell ; correction et évaluation des rendus.
	\item Création de {\bfseries contenus pédagogiques interactifs} (quiz, cours inversés, tutoriels).
\end{itemize}

\par\vspace{5pt}

{\large\bfseries Projets Personnels}
\par\vspace{-8pt}
\rule{\linewidth}{0.4pt}
\par\vspace{5pt}

\underline{\bfseries Sledge -- Moteur de migration, synchronisation et orchestration multi-cloud}
\par
\begin{adjustwidth}{5pt}{0pt}
    Outil générique de migration multi-cloud gérant des structures de données complexes et des bases de données. \\
    Architecture et livraison d'une plateforme backend complète en Go, optimisée pour des flux de données à haut volume, via une couche d'abstraction cloud personnalisée ({\ttfamily provicol}) supportant AWS S3 et GCP. \\
    Utilisation de l'API Docker à l'exécution pour orchestrer des conteneurs d'environnements concurrents.
\end{adjustwidth}
\par
{\bfseries Compétences :} Go, API Docker, Python, SQL, intégrité des données, automatisation système

\par\vspace{8pt}

\underline{\bfseries Area -- Plateforme d'automatisation de services (clone d'IFTTT)}
\par
\begin{adjustwidth}{5pt}{0pt}
    Conception et développement seul du backend en Go avec PostgreSQL, en architecture microservices (services en Go et Python), derrière un proxy Traefik. \\
    Intégration d'une quinzaine de services tiers via OAuth, webhooks et requêtes HTTP.
\end{adjustwidth}
\par
{\bfseries Compétences :} Go, Python, PostgreSQL, microservices, Traefik, OAuth, webhooks

\par\vspace{8pt}

\underline{\bfseries Chiffrement -- Implémentations AES et RSA}
\par
\begin{adjustwidth}{5pt}{0pt}
    Implémentation from scratch des algorithmes {\bfseries AES et RSA} ({\ttfamily myAES}, {\ttfamily myRSA}) pour la compréhension des mécanismes cryptographiques utilisés dans les blockchains et les certificats numériques.
\end{adjustwidth}
\par
{\bfseries Compétences :} C, cryptographie, AES, RSA

\par\vspace{8pt}

\underline{\bfseries Zappy -- Serveur TCP multi-clients en C}
\par
\begin{adjustwidth}{5pt}{0pt}
    Serveur TCP mono-thread à boucle événementielle ({\bfseries \texttt{select}}), avec intégration Lua en C ; aucune erreur ni fuite sous Valgrind.
\end{adjustwidth}
\par
{\bfseries Compétences :} C, sockets TCP, select, Lua, gestion mémoire, Valgrind

\par\vspace{8pt}

{\large\bfseries Formation}
\par\vspace{-8pt}
\rule{\linewidth}{0.4pt}
\par\vspace{5pt}

{\bfseries Epitech, Paris} \hfill 2022 -- 2027 \\
Diplôme d'ingénieur en technologies de l'information (diplôme de niveau Master) \\
Certification Niveau 7 - Expert en ingénierie logicielle

\vspace{6pt}

{\bfseries McGill University, Montréal} \hfill 2025 -- 2026 \\
Certificat en Management -- GPA : 3.57/4.0 \\
Cours suivis : Comptabilité, Managerial Economics \& Analysis, Business, Mathématiques, Gestion de projet, Marketing

\end{document}
